%%%%%
%
%% Dans ce fichier, les �ventuels accents sont cod�s
%% en ISO-8859-15 (aka latin9)
%% La ligne suivante donne trois exemples 
%% e accent aigu : �, e accent grave : �, e-dans-l'o : �
%
\section{Point de ros�e ou de givre sachant l'humidit� relative}

\itemstyle{Op�rateur~:} \operfmt{Tdf\_Hu}
\index{Tdf\_Hu}

%%%%%OBJET%%%%
\itemstyle{Objet~:}
 Calcul de la temp�rature du point de ros�e ou de givre � partir
de l'humidit� relative � l'eau liquide et la temp�rature de l'air.

%%%%%ARGUMENTS%%%%
\bigskip
%% Dans ce fichier, les �ventuels accents sont cod�s
%% en ISO-8859-15 (aka latin9)
%% La ligne suivante donne trois exemples 
%% e accent aigu : �, e accent grave : �, e-dans-l'o : �
%%%%%ENTREES%%%%
\begin{tabarg}
\multicolumn{3}{l}{\titretab {Entr\'ees}}\\
\hline
$Hu$ & Vecteur, \texttt{float} & Humidit� relative (\%)\\
$Ts$ & Vecteur, \texttt{float} & Temp�rature de l'air (K ou Celsius)\\
\end{tabarg}

\medskip
%%%%%SORTIES%%%%
\begin{tabarg}
\multicolumn{3}{l}{\titretab {Sorties}}\\
\hline
$Tdf$ & Vecteur, \texttt{float} & Temp�rature du point de ros�e ou de givre (Celsius)\\
\end{tabarg}


%%%%%FORMULE%%%%
\itemstyle{Formule}

%% Formule de l'operateur : Tdf_Hu
%% Dans ce fichier, les �ventuels accents sont cod�s
%% en ISO-8859-15 (aka latin9)
%% La ligne suivante donne trois exemples 
%% e accent aigu : �, e accent grave : �, e-dans-l'o : �
$Ts$ est converti en Kelvin, si besoin est, puis :
\[
e = \frac{Hu}{100} . e_w(Ts) \qquad \text{(d�finition de Hu)}
\]
et
\[
Tdf_K = 
\begin{cases}
e_{w}^{-1}(e) \qquad \text{ si } e > 6.11 hPa \\
e_{i}^{-1}(e) \qquad \text{ sinon }  \\
\end{cases}
\]
enfin, $Tdf = Tdf_K -273,15 $.


%%%%%AUTEUR%%%%
\itemstyle{Auteur~:}
CNRM/GMEI/TRAMM

%%%%%REFERENCE%%%%
\itemstyle{R\'ef\'erences~:}

Les formules directes et r�ciproques de $ew(T)$ sont
celles de l'article Bolton-1980 \cite{Bolton-80} :
eq. 10, page 1047, d�crite en annexe 
\ref{annexe-ewbolton}.

La formule r�ciproque de $ei(T)$ est celle 
recommand�e par l'OMM dans l'annexe 4.B du  \emph{Guide des instruments
et des m�thodes d'observation m�t�orologique} \cite{OMM-N8-1996},
d�crite en annexe \ref{annexe-invewiOMM} de ce propre document.
